%% -*- coding:utf-8 -*-

\chapter{Der Verbalkomplex}
\label{chap-verbal-complex}


\section{Das Phänomen}

SOV"=Sprachen\is{Verbalkomplex|(} wie das \ili{Niederländisch}e und das \ili{Deutsch}e bilden Verbalkomplexe, das heißt
Kombinationen von Verben unter Ausschluss der nicht"=verbalen Argumente der Verben. So nimmt man zum
Beispiel an, dass \emph{zu lesen}, \emph{versprochen} und \emph{hat} in Haiders Beispiel (\citeyear[\page
  110]{Haider86c}; \citeyear[\page 128]{Haider90b}) in (\mex{1}) unter Ausschluss der Argumente der Verben
\emph{es}, \emph{ihr} und \emph{jemand} eine Konstituente bilden.
\ea\label{ex-weil-es-ihr-jemand-zu-lesen-versprochen-hat}
weil es ihr jemand zu lesen versprochen hat
\z
Es gibt mehrere Indikatoren für die Bildung von Verbalkomplexen,
die von Gunnar \citet{Bech55a} im Detail herausgearbeitet wurden. Wie oben angedeutet, besteht eine Möglichkeit, solche
Verbalkomplexe zu analysieren, darin anzunehmen, dass die Verben in einem Satz eine Einheit bilden, die sich im Wesentlichen wie ein
einzelnes Verb verhält. Das erklärt zum Beispiel, warum die Argumente der drei Verben in (\mex{0}) gescrambelt werden können:
\emph{es} hängt von \emph{zu lesen} ab, \emph{ihr} hängt von \emph{versprochen}
ab und \emph{jemand} ist das Subjekt und kongruiert mit dem finiten Verb \emph{hat}
(normalerweise wird es auch als ein Abhängiges des Auxiliars \emph{hat} behandelt).

Es sei angemerkt, dass es in den deutschen\il{Deutsch} Dialekten extreme Variation gibt, was die Serialisierung
der Elemente im Verbalkomplex betrifft. Das regierende\footnote{
  Der Begriff \emph{regieren}\is{Rektion} wird gleichbedeutend mit \emph{selegieren} verwendet. Wenn ein Verb ein Akkusativ"=
  objekt verlangt, sagt man, dass es den Akkusativ/ein Akkusativobjekt regiert. Verben können auch andere Verben regieren
  und bestimmen, welche Form das regierte Verb annehmen muss, \zb Perfekt-/Passivpartizip, Infinitiv
  mit oder ohne \emph{zu}.
} Verb wird im Standard\-deutschen\il{Deutsch} rechts vom
eingebetteten Verb realisiert: V$_3$ V$_2$ V$_1$ wie in (\mex{0}), aber es gibt Beispiele wie
(\mex{1}) aus \citew[376]{Mueller99a}.\footnote{
Interviewpartner in: \emph{Insekten und andere Nachbarn -- ein Haus in Berlin}, ARD 1995-11-15.
}
\eal
\ex
Ich hätte stapelweise Akten kön\-nen haben.\hfill (deutscher\il{Deutsch} Berliner Dialekt)
\ex
weil ich mir das  nich hab' lassen gefallen
\ex
wenn se   mir hier würden rausschmeißen, \ldots
\zl
Die Abfolgen in (\mex{0}) entsprechen der Abfolge, die im \ili{Niederländisch}en am natürlichsten ist. (\mex{1}) zeigt einige
niederländische\il{Niederländisch} Beispiele:
\eal
\ex
\gll dat   Kim het boek wil lezen\\
     dass  Kim das Buch will lesen\\
\glt `dass Kim das Buch lesen will'
\ex
\gll dat  Kim Sandy het boek laat lezen\\
     dass Kim Sandy das Buch lässt lesen\\
\glt `dass Kim Sandy das Buch lesen lässt'
\ex
\gll dat Kim Sandy het boek wil laten lezen\\
     dass Kim Sandy das Buch will lassen lesen\\
\glt `dass Kim Sandy das Buch lesen lassen will'
\zl

SVO"=Sprachen wie das \ili{Dänisch} und das \ili{Englisch} erlauben es nicht, dass die Argumente eingebetteter Verben
mit Argumenten höherer Verben gescrambelt werden. Alle Argumente bleiben in ihrer VP (natürlich modulo Extraktion).


\section{Die Analyse}


Die Technik, die zur Analyse der Verbalkomplexe verwendet wird, heißt \emph{\isi{Argumentattraktion}} oder
\emph{Argumentkomposition} und wurde von \citet{Geach70a} im Rahmen der Kategorialgrammatik
entwickelt und von \citet{HN94a} für die HPSG adaptiert. Die Analyse von \emph{lesen wird}, wie es in
(\mex{1}) vorkommt, ist in Abbildung~\vref{fig-lesen-wird} dargestellt.
\ea
dass keiner das Buch lesen wird
\z
\begin{figure}
\begin{forest}
sm edges
[V\feattab{
%              \vform \type{fin},\\
              \sliste{ NP[\type{nom}], NP[\type{acc}] } } 
        [V\feattab{
%              \vform \type{bse},\\
              \sliste{ NP[\type{nom}], NP[\type{acc}]} }, name=lesen [lesen] ]
        [V\feattab{
%              \vform \type{fin},\\
              \sliste{ NP[\type{nom}], NP[\type{acc}], V }}, name=wird [wird] ]
]
\draw[semithick,->] (lesen)..controls +(south east:2) and +(south west:2)..(wird);
\end{forest}
\caption{\label{fig-lesen-wird}Analyse der Verbalkomplexbildung von \emph{lesen wird}
  mittels Argumentkomposition (vorläufige Version)}
\end{figure}
\emph{wird} selegiert einen Infinitiv ohne \emph{zu} und zusätzlich dessen Argumente. Dieser
Infinitiv (\emph{lesen}) wird mit dem Verb kombiniert und ist daher nicht in der Valenzliste
des Mutterknotens enthalten.

Um auf unsere Einkaufsanalogie von S.\,\pageref{page-shopping-analogy} zurückzukommen: Die Verbalkomplexbildung
kann man sich vorstellen, indem man sich ein junges und hilfsbereites Hilfsverb vorstellt, das einer Person aus
einer Risikogruppe mitten in einer Pandemie aushilft. Da Risikopersonen nicht einkaufen gehen
sollen, übernimmt die hilfsbereite Person deren Einkaufsliste und erledigt den Einkauf für sie. Im
Fall von Hilfsverben selegiert das Hilfsverb einfach das Hauptverb und
verlangt keine weiteren Argumente außer denen, die vom eingebetteten Verb übernommen werden. Das heißt, das
Hilfsverb erledigt einfach den Einkauf für das Hauptverb. Ein sehr altruistisches Verb ist es. Später werden wir
uns Verben wie \emph{versuchen} und \emph{lassen} anschauen, die zusätzlich zu denen des eingebetteten Verbs eigene Argumente
verlangen. Das ist parallel zu einem Einkaufsereignis, bei dem die helfende Person
zusätzlich zum Einkauf für jemand anderen ihre eigenen Waren kauft.

%Now, back to linguistics and to some details of the analysis:
Die Kombination von \emph{lesen} und \emph{wird} verhält sich wie ein einzelnes Verb, indem sie
in beliebiger Reihenfolge mit ihren Argumenten kombiniert werden kann. Abbildung~\ref{fig-vc-nom-acc} zeigt die Analyse von (\mex{1}a) und
Abbildung~\ref{fig-vc-acc-nom} zeigt die Analyse von (\mex{1}b).
\eal
\ex\label{ex-dass-keiner-das-buch-lesen-wird}
{[}dass{]} keiner das Buch lesen wird
\ex
{[}dass{]} das Buch keiner lesen wird
\zl

\begin{figure}
\centerfit{
\begin{forest}
sm edges
[V\feattab{
              \sliste{ }}
        [{NP[\type{nom}]} [keiner] ]
        [V\feattab{
              \sliste{ NP[\type{nom}] }}, s sep+=1em
          [{NP[\type{acc}]} [das Buch, roof] ]
          [V\feattab{
%              \vform \type{fin},\\
              \sliste{ NP[\type{nom}], NP[\type{acc}]}}
             [V\feattab{
%              \vform \type{bse},\\
              \sliste{ NP[\type{nom}], NP[\type{acc}]}} [lesen] ]
             [V\feattab{
%              \vform \type{fin},\\
                \sliste{ NP[\type{nom}], NP[\type{acc}], V }} [wird] ] ] ] ]
\end{forest}}
\caption{\label{fig-vc-nom-acc}Bildung eines Verbalkomplexes und Realisierung der Argumente in normaler
  Abfolge (vorläufige Version)}
\end{figure}




\begin{figure}
\centerfit{
\begin{forest}
sm edges
[V\feattab{
              \sliste{ }}
        [{NP[\type{acc}]} [das Buch, roof] ]
        [V\feattab{
              \sliste{ NP[\type{acc}] }}
          [{NP[\type{nom}]} [keiner] ]
          [V\feattab{
%              \vform \type{fin},\\
              \sliste{ NP[\type{nom}], NP[\type{acc}] }}
             [V\feattab{
%              \vform \type{bse},\\
              \sliste{ NP[\type{nom}], NP[\type{acc}]}} [lesen] ]
             [V\feattab{
%              \vform \type{fin},\\
                \sliste{ NP[\type{nom}], NP[\type{acc}], V }} [wird] ] ] ] ]
\end{forest}}
\caption{\label{fig-vc-acc-nom}Bildung eines Verbalkomplexes und Scrambling der Argumente (vorläufige Version)}
\end{figure}

\noindent
Ich folge
\citet[Abschnitt~3.1.1]{Kiss95a} und repräsentiere das Subjekt nicht"=finiter Verben als den Wert eines
speziellen Merkmals \subj. \subj unterscheidet sich von \spr und \comps darin, dass es kein Valenzmerkmal ist. Der
Grund für diese Sonderbehandlung ist, dass das Subjekt nicht als Teil einer nicht"=finiten
Verbalphrase realisiert werden kann. Das ist besonders klar bei Infinitiven mit \emph{zu}:
\eal
\ex[]{
Aicke hat Conny versprochen, [das Buch zu lesen].
}
\ex[*]{
Aicke hat Conny versprochen, [sie das Buch zu lesen].
}
\zl
Andere nicht"=finite Verben (reine Infinitive und Partizipien) können nicht ins \nf{} gestellt werden, aber sie
können vorangestellt werden. (\mex{1}) zeigt, dass das Subjekt nicht zusammen mit anderen Argumenten
des Verbs ins \vf{} gestellt werden kann.
\eal
\judgewidth{?*}
\ex[]{
{[}Das Buch lesen{]} wird Aicke morgen.
}
\ex[*]{
{[}Aicke lesen{]} wird das Buch morgen.
}
\ex[?*]{
{[}Aicke das Buch lesen{]} wird morgen.
}
\zl
Der Lexikoneintrag für die nicht"=finite Form von \emph{lesen} ist in (\mex{1}) angegeben:
\ea
\emph{lesen} nicht"=finite Form:\\
\ms{
subj  & \sliste{ NP[\type{nom}] }\\
comps & \sliste{ NP[\type{acc}] }\\
}
\z

\noindent
Die folgende Merkmal-Wert-Matrix (AVM)\is{Merkmal-Wert-Matrix (AVM)} ist eine Repräsentation des Hilfsverbs \emph{werden}:
\eas
\emph{werden} nicht"=finite Form:\\
\ms{
subj  & \ibox{1}\\
comps & \ibox{2} $\oplus$ \sliste{ V[\vform \type{bse}, \textsc{lex}+, \subj \ibox{1}, \comps \ibox{2}] }\\
}
\zs
\emph{werden} selegiert ein Verb, das die \type{bse}-Form hat, also einen Infinitiv ohne
\emph{zu}. Das eingebettete Element muss lexikalisch sein (\textsc{lex}+), das heißt ein einzelnes Wort oder ein
Verbalkomplex. Von allen Phrasen, die durch das Kopf-Komplement-Schema und das Spezifikator-Kopf-Schema
lizenziert werden, wird angenommen, dass sie \textsc{lex}$-$ sind.
Die Kästchen mit Zahlen sind im Wesentlichen Variablen. Ihre Werte hängen von den Werten der
eingebetteten Verben ab. Daher kann dieser Lexikoneintrag mit einem Verb wie \emph{lesen}
verwendet werden, das ein nominativisches und ein akkusativisches Argument nimmt, aber auch mit einem Verb wie \emph{helfen},
das ein nominativisches und ein dativisches Argument nimmt.

Bevor ich auf die Details der Analyse eingehe, muss ich die Lexikoneinträge für die finite
Form der Auxiliare bereitstellen. Da das Subjekt finiter Verben natürlich realisiert werden kann, muss es
in einer der Valenzlisten repräsentiert werden. Wie in Abschnitt~\ref{sec-intro-spr-comps} besprochen,
werden deutsche\il{Deutsch} Subjekte in der \comps-Liste finiter Verben repräsentiert. Daher hat der Lexikoneintrag für
\emph{wird} die folgende Form:
\eas
\emph{wird} finite Form:\\
\ms{
subj  & \sliste{}\\
comps & \ibox{1} $\oplus$ \ibox{2} $\oplus$ \sliste{ V[\vform \type{bse}, \textsc{lex}+, \subj \ibox{1}, \comps \ibox{2}] }\\
}
\zs
%\largerpage
Das besagt im Wesentlichen, dass die Valenz von \emph{wird} aus einem eingebetteten Verb besteht und aus dem, was die
\subjl{} dieses Verbs ist, plus dem, was die \comps-Liste dieses Verbs ist. Das wird für
\emph{lesen wird} in Abbildung~\vref{fig-lesen-wird-details} veranschaulicht.\footnote{%
Die Lexikoneinträge komplexbildender Prädikate verlangen, dass ihr verbales Argument \lex{}$+$ ist. So ist in
Abbildung~\ref{fig-lesen-wird-details} auch \emph{lesen} \lex{}$+$. Da diese Information
für die Diskussion der Argumentattraktion nicht relevant ist, wird sie in
Abbildung~\ref{fig-lesen-wird-details} und den folgenden Abbildungen weggelassen. \emph{lesen} und \emph{lesen können}
müssen in den Verbalkomplexen, die in den
Abbildungen~\ref{fig-lesen-koennen-wird} und~\ref{fig-wird-lesen-koennen} dargestellt sind, \lex{}$+$ sein. Natürlich sind Lexikoneinträge wie
\emph{wird} und \emph{können} ebenfalls \lex{}$+$, weil sie
Wörter sind. Diese Information in den Bäumen darzustellen, wäre eher verwirrend als erklärend, und daher
habe ich beschlossen, die \lex-Information wegzulassen.
}
\begin{figure}
\begin{forest}
sm edges
[V\feattab{
              \vform \type{fin},\\
              \comps \ibox{1} $\oplus$ \ibox{2} } 
        [{\ibox{3} V}\feattab{
              \vform \type{bse},\\
              \subj  \ibox{1} \sliste{ NP[\type{nom}] }, \\ 
              \comps \ibox{2} \sliste{ NP[\type{acc}] } } [lesen] ]
        [V\feattab{
              \vform \type{fin},\\
              \comps \ibox{1} $\oplus$ \ibox{2} $\oplus$ \sliste{ \ibox{3} } } [wird] ] ]
\end{forest}
\caption{\label{fig-lesen-wird-details}Detaillierte Analyse eines Verbalkomplexes}
\end{figure}
Das Auxiliar selegiert einen Infinitiv ohne \emph{zu} \iboxb{3}. Das wird durch den Wert
\type{bse} für die \vformf{} des selegierten Verbs sichergestellt: \type{bse}\istype{bse} steht für einen Infinitiv ohne
\emph{to}/\emph{zu}/\ldots{}, \type{inf}\istype{inf} steht für eine Infinitivform mit Marker, \type{ppp}\istype{ppp}
steht für ein Partizip und \type{fin}\istype{fin} für ein finites Verb. Das Subjekt des selegierten Infinitivs
\iboxb{1} und die Komplemente \iboxb{2} werden übernommen. Das Ergebnis ist, dass \emph{lesen wird} dieselben
Argumente hat wie \emph{liest}.

Um das Ganze noch unterhaltsamer zu machen, können wir es komplexer gestalten und uns Verbalkomplexe mit
drei Verben anschauen. Abbildung~\vref{fig-lesen-koennen-wird} zeigt die Analyse des Verbalkomplexes \emph{lesen können wird}
in Sätzen wie (\mex{1}):
\ea
\label{ex-lesen-koennen-wird}
{[}dass{]} sie das Buch lesen können wird
\z

\noindent
Die Analyse von (\mex{0}) ist in Abbildung~\ref{fig-lesen-koennen-wird} angegeben.
\begin{figure}
\centerfit{%
\begin{forest}
sm edges
[V\feattab{
              \vform \type{fin},\\
              \comps \ibox{1} $\oplus$ \ibox{2} } 
        [{\ibox{4} V\feattab{
              \vform \type{bse},\\
              \subj  \ibox{1},\\
              \comps \ibox{2} }} 
           [{\ibox{3} V\feattab{
              \vform \type{bse},\\
              \subj  \ibox{1} \sliste{ NP[\type{nom}] }, \\ 
              \comps \ibox{2} \sliste{ NP[\type{acc}] } }} [lesen] ]
           [V\feattab{
              \vform \type{bse},\\
              \subj  \ibox{1},\\
              \comps \ibox{2} $\oplus$ \sliste{
                \ibox{3} } } [können] ] ]
        [V\feattab{
              \vform \type{fin},\\
              \comps \ibox{1} $\oplus$ \ibox{2} $\oplus$ \sliste{
                \ibox{4} } } [wird] ]
]
\end{forest}}
\caption{\label{fig-lesen-koennen-wird}Analyse eines deutschen Verbalkomplexes mit drei Verben in kanonischer Abfolge}
\end{figure}

%%\largerpage[2]
Ein interessanter Aspekt der Analyse ist, dass sie ein Phänomen erklären kann, das Auxiliary
Flip\is{Auxiliary Flip} oder \emph{Oberfeldumstellung} genannt wird. Das \ili{Deutsch}e erlaubt es optional, Verben, die ein Modalverb\is{Verb!Modal} regieren,
links vom Verbalkomplex statt rechts vom Modalverb zu platzieren. Statt
(\ref{ex-lesen-koennen-wird}) kann man also auch die Abfolge in (\mex{1}) verwenden:
\ea
{[}dass{]} sie das Buch wird lesen können
\z
\begin{figure}
\centerfit{%
\begin{forest}
sm edges
[V\feattab{
              \vform \type{fin},\\
              \comps \ibox{1} $\oplus$ \ibox{2} } 
        [V\feattab{
              \vform \type{fin},\\
              \comps \ibox{1} $\oplus$ \ibox{2} $\oplus$ \sliste{
                \ibox{4} } } [wird] ]
        [{\ibox{4} V\feattab{
              \vform \type{bse},\\
              \subj  \ibox{1},\\
              \comps \ibox{2} }} 
           [{\ibox{3} V\feattab{
              \vform \type{bse},\\
              \subj  \ibox{1} \sliste{ NP[\type{nom}] }, \\ 
              \comps \ibox{2} \sliste{ NP[\type{acc}] } }} [lesen] ]
           [V\feattab{
              \vform \type{bse},\\
              \subj  \ibox{1},\\
              \comps \ibox{2} $\oplus$ \sliste{
                \ibox{3} } } [können] ] ]
]
\end{forest}}
\caption{\label{fig-wird-lesen-koennen}Analyse eines deutschen Verbalkomplexes mit drei Verben mit
  Auxiliary Flip}
\end{figure}

\noindent
Nachdem ich die Analyse von Verbalkomplexen besprochen habe, wie sie aus den OV-Sprachen wie dem
\ili{Deutsch}en, \ili{Niederländisch}en und \ili{Afrikaans} bekannt sind, möchte ich kurz auf die SVO-Sprachen wie das \ili{Dänisch} und das \ili{Englisch} eingehen. Normalerweise verlangt ein Kopf, dass sein Argument gesättigt ist, das heißt, die
\compsv{} muss bei NP, PPs, APs, CPs und satzwertigen sowie VP-Argumenten die leere Liste sein. Verbalkomplexe sind anders: Wörter werden direkt kombiniert. Die
VO-Sprachen unterscheiden sich von den OV-Sprachen darin, dass sie dies nicht erlauben. In VO-Sprachen bildet das Verb eine
Phrase mit seinen Komplementen, und diese Verbalphrase kann unter ein anderes Verb eingebettet werden. (\mex{1}a)
zeigt ein Beispiel mit Hilfsverben, (\mex{1}b) ist ein Beispiel mit einem Vollverb, das zusätzlich eine
Infinitiv-Verbalphrase mit \emph{to} und ein Objekt nimmt.
\eal
\ex
\gll Kim [will [have [read the book]]].\\
     Kim {[}wird {[}haben {[}gelesen das Buch{]}{]}{]}\\
\glt `Kim wird das Buch gelesen haben.'
\ex
\gll Somebody [promised her [to read it]].\\
     jemand {[}versprach ihr {[}zu lesen es{]}{]}\\
\glt `Jemand hat ihr versprochen, es zu lesen.'
\zl
Sprachen wie das \ili{Dänisch} und das \ili{Englisch} haben nur das Kopf-Komplement-Schema und das Spezifikator-Kopf-
Schema, während Sprachen wie das \ili{Niederländisch} und das \ili{Deutsch} ein zusätzliches Schema haben, das ungesättigte
Wörter kombinieren kann.
%\todostefan{check whether this has to be explained better}
Das Schema für die Prädikatkomplexbildung ist in Abbildung~\vref{fig-pred-complex} skizziert.
\begin{figure}
\begin{forest}
[{[\comps \ibox{1}]}
  [\ibox{2} ]
  [{[\comps \ibox{1} $\oplus$ \sliste{ \ibox{2} }]}]]
\end{forest}
\caption{\label{fig-pred-complex}Skizze des Prädikatkomplex"=Schemas\is{Schema!Prädikatkomplex}}
\end{figure}
Dieses Schema ist dem Kopf-Komplement-Schema sehr ähnlich, das auf
Seite~\pageref{fig-head-comp} angegeben wurde. Der Unterschied besteht darin, dass dieses Schema keine \lex{}$-$-Elemente
lizenziert, wie es das Spezifikator-Kopf- und das Kopf-Komplement-Schema tun. Daher ist die Kombination zweier Verben
mit \lex{}+-Anforderungen regierender Verben kompatibel, und eine Einbettung in noch komplexere Verbal\-
komplexe ist möglich. Abbildung~\ref{fig-lesen-koennen-wird} ist ein Beispiel: Die Kombination von
\emph{lesen} und \emph{können} ist mit der \lex{}+-Anforderung von \emph{wird}
kompatibel. Das Prädikatkomplex-Schema unterscheidet sich vom Kopf-Komplement-Schema im \ili{Deutsch}en auch darin, dass
es den Kopf mit dem letzten Element der Valenzliste kombiniert. Das ist das eingebettete Verb, und es muss
vor jedem anderen Argument mit dem Kopf kombiniert werden, da man nicht wüsste, was die anderen
Argumente sind, weil sie vom eingebetteten Verb übernommen werden.

Bevor ich mich dem nächsten Phänomen zuwende, möchte ich kurz eine Alternative zu der hier vorgestellten
Verbalkomplexanalyse besprechen. Ein alternativer Vorschlag bestand darin, Auxiliare im \ili{Deutsch}en als VP-
einbettende Verben zu analysieren \citep{Wurmbrand2003b}. Unser Standardbeispiel hätte dann die Analyse in (\mex{1}):
\ea
dass keiner [[das Buch lesen] wird]
\z
Die Frage, die solche Analysen beantworten müssen, ist, wie das Scrambling der Argumente der beteiligten Verben
erklärt werden kann. Die Antwort ist oft, dass man annimmt, dass das Objekt des eingebetteten Verbs
aus der VP extrahiert und an die linke Peripherie des Satzes bewegt wird. Das wird in (\mex{1}) gezeigt:
\ea
dass [das Buch]$_i$ keiner [[ \_$_i$ lesen] wird]
\z
Analysen, die Scrambling als Bewegung behandeln, sind jedoch problematisch, da sie zusätzliche
Lesarten von Sätzen vorhersagen, die Quantoren in ihren NPs haben (\citealp[\page 146]{Kiss2001a};
\citealp[Abschnitt~2.6]{Fanselow2001a}).\todostefan{Change discussion}


Bevor ich mich der Analyse der Verbstellung zuwende, möchte ich zeigen, wie Sätze mit mehreren Verben
\label{page-English-Aux-VPs}
in SVO-Sprachen analysiert werden können. Abbildung~\vref{fig-vp-embedding-svo} zeigt die Analyse der englischen\il{Englisch} Version von Satz
(\ref{ex-dass-keiner-das-buch-lesen-wird}).
\begin{figure}
\centerfit{
\begin{forest}
sm edges
[S
        [{NP[\type{nom}]} [nobody] ]
        [VP
          [V%% \feattab{
            %%   spr \sliste{ NP[\type{nom}] }\\
            %%   comps \sliste{ VP }} 
              [will] ]
          [VP%% \feattab{
             %%  spr \sliste{ NP[\type{nom}] }}
            [V%% \feattab{
              %% spr \sliste{ NP[\type{nom}] }\\
              %% comps \sliste{ NP[\type{acc}] }} 
              [read] ]
            [{NP[\type{acc}]} [the book, roof] ] ] ] ]
\end{forest}}
\caption{\label{fig-vp-embedding-svo}Einbettung einer VP in SVO-Sprachen}
\end{figure}
Das Verb \emph{reads} selegiert ein Subjekt und ein Objekt. Das Verb bildet eine VP mit der NP \emph{the
  book}. Dieser VP fehlt noch ein Subjekt. Das Auxiliar \emph{will} selegiert eine VP und ein Subjekt,
das mit dem Subjekt von \emph{read} identisch ist. Die Kombination von \emph{will} und der VP wird
durch das Kopf-Komplement-Schema lizenziert, das in Abbildung~\vref{fig-head-comp} skizziert wurde.

Das Äquivalent von \emph{lesen können wird} kann hier nicht angegeben werden, da englische\il{Englisch} Modalverben\is{Verb!Modal}
keine nicht"=finiten Formen haben, aber man kann Beispiele mit Modalverben als höchstem Verb konstruieren:
\ea
\gll She [must [have [seen it]]].\\
     sie {[}muss {[}haben {[}gesehen es{]}{]}{]}\\
\glt `Sie muss es gesehen haben.'
\z
Dieser Satz hat eine Struktur, die der in Abbildung~\ref{fig-vp-embedding-svo} ähnlich ist:
\emph{must} und \emph{have} betten beide VPs ein.

Schließlich zeigt Abbildung~\vref{fig-somebody-has-promised-him-to-read-it} die Übersetzung von
(\ref{ex-weil-es-ihr-jemand-zu-lesen-versprochen-hat}):
\ea
\gll Somebody has promised her to read it.\\
     jemand hat versprochen ihr zu lesen es\\
\glt `Jemand hat ihr versprochen, es zu lesen.'
\z
\emph{promise} ist ein Verb, das ein Subjekt, ein Objekt und ein VP-Komplement nimmt.
\begin{figure}
\centerfit{
\begin{forest}
sm edges
[S
        [{NP[\type{nom}]} [somebody] ]
        [VP
            [V [has]]
            [VP 
              [\vbar [V [promised] ]
                     [{NP[\type{acc}]} [her] ] ]
              [VP [V [to]]
                [VP
                  [V [read] ]
                  [{NP[\type{acc}]} [the book, roof] ] ] ] ] ] ]
\end{forest}}
\caption{\label{fig-somebody-has-promised-him-to-read-it}Einbettung einer VP mit Verben, die ein
  zusätzliches Objekt nehmen}
\end{figure}
Wie in der Analyse von (\ref{ex-nobody-gave-the-child-a-book}) auf
Seite~\pageref{ex-nobody-gave-the-child-a-book} -- die hier der Bequemlichkeit halber als (\mex{1}) wiederholt wird
-- wird das Verb \emph{promised} zuerst mit seinem NP-Komplement und dann mit seinem VP-
Argument kombiniert.
\ea
\label{ex-nobody-gave-the-child-a-book-three}
\gll Nobody gave the child a book.\\
     niemand gab dem Kind ein Buch\\
\glt `Niemand gab dem Kind ein Buch.'
\z
Das VP-Argument von \emph{promised} in Abbildung~\ref{fig-somebody-has-promised-him-to-read-it} besteht aus \emph{to} und einer weiteren VP mit einem Infinitiv in Grundform. \emph{to} wird
als Hilfsverb analysiert \parencites[\page 600]{GPS82a-u}[\page 147]{Sag2020a}.
Es ist wichtig anzumerken, dass das Objekt \emph{him} in keiner anderen Position erscheinen kann (abgesehen von
einer Extraktion an die linke Peripherie). Zum Beispiel kann es nicht in der Position von \emph{the book} erscheinen,
und dasselbe gilt für \emph{the book}: Diese Phrase kann an keiner anderen Stelle als in der
Objektposition erscheinen.
\is{Verbalkomplex|)}

\exercises{


\begin{enumerate}
\item Skizzieren Sie die Analyse der Verbalkomplexe in den folgenden Beispielen:
\eal
\ex
dass sie darüber lachen muss
\ex
dass sie darüber hat lachen müssen
\ex
dass sie darüber     wird haben lachen müssen
\zl
Geben Sie bitte die \subj- und die \comps-Werte an. Die \spr-Werte können Sie weglassen, da sie für alle deutschen\il{Deutsch} Verben ohnehin die leere Liste sind.


\item Suchen Sie zwei Sätze mit einem Verbalkomplex in einer Zeitung oder in Korpora (zum Beispiel im
  COSMAS-Korpus\footnote{%
\url{https://cosmas2.ids-mannheim.de}, 2020-05-11.}) und
  analysieren Sie die Verbalkomplexe.

\item Suchen Sie in einem Korpus nach Verbalkomplexen mit mehr als vier Verben und dokumentieren Sie Ihre Suche.
\end{enumerate}

}

\furtherreading{Die Analyse von Verbalkomplexen in der HPSG wurde zuerst von \citet{HN89b,HN89a} entwickelt.
\citew{HN94a} ist die erste begutachtete Publikation zu diesem Thema von Hinrichs \&
Nakazawa. \citew{Kiss95a} ist eine Monographie, die sich mit Verbalkomplexen befasst. \citet{Meurers2000b} befasst sich
ebenfalls mit Verbalkomplexen einschließlich schwieriger Fälle wie der sogenannten \emph{Zwischenstellung},
die hier nicht behandelt wird. \citet{Mueller2002b} behandelt nicht nur Verbalkomplexe, sondern auch andere Typen
komplexer Prädikate wie Adjektiv-Verb-Komplexe, Resultativkonstruktionen und Partikelverben.

\citet{GS2021a} geben einen Überblick über Analysen komplexer Prädikate in der HPSG.

\citet[Kapitel~7]{Haider2010a} diskutiert verschiedene Vorschläge im Rahmen der Government \& Binding.}



%      <!-- Local IspellDict: de_DE -->
